\subsection{Two objects can carry nontrivial deck coupling}

\begin{lemma}[Representing a corner map by a finite biset]
\label{D:corner-biset}
Let $P,M$ be finite monoids and let $\phi:P\to eMe$ be a unital
homomorphism, where $e\in M$ is idempotent. Regard $Me$ as an $(M,P)$-biset
using $\phi$ on the right. Its category of elements has classifying space
homotopy equivalent to $BP$; under this equivalence its two projections
induce the identity of $BP$ and the corner map $BP\to BM$ defined by $\phi$.
Consequently, for any finite family $\phi_i:P\to e_iMe_i$, the two-object
category with endomorphism monoids $P,M$, no arrows in the reverse
direction, and cross-biset $\coprod_i Me_i$ realizes
\[
 \operatorname{hocolim}\left(
 BP\ \longleftarrow\ \coprod_i BP\ \longrightarrow\ BM\right),
\]
with left maps the identity and right maps the prescribed corner maps.
\end{lemma}
\begin{proof}
For one component, the category of elements has objects $x\in Me$ and
arrows $x\to x'$ given by pairs $(p,m)\in P\times M$ satisfying
$x'\phi(p)=mx$. Projection to $P$ is a Grothendieck fibration:
the cartesian lift over $p$ with target $x'$ has source $x'\phi(p)$
and second component $1$. Its fiber is the action category $M\ltimes Me$.
The $M$-set $Me$ is a retract of the regular $M$-set $M$, by
$m\mapsto me$ and inclusion. The action category $M\ltimes M$ has initial
object $1$. Therefore $B(M\ltimes Me)$ is a retract of a contractible
space and is contractible. The fiber-to-comma adjunction for the displayed
fibration, or Quillen's Theorem A applied to these cartesian lifts, makes
projection to $BP$ a homotopy equivalence.

In the Karoubi category of $M$, the arrows $x:e\to1$ give a natural
transformation between the two projected functors: its naturality
identity is exactly $x'\phi(p)=mx$. This identifies the other projection
with the specified corner map. The last assertion is the collage
homotopy-pushout lemma checked in Lemma~\ref{lem:B-prelim}.
\end{proof}

\begin{proposition}[A two-object degree-two mapping torus]
\label{D:degree-two-torus}
The mapping torus of a degree-two self-map of $S^2$ is the classifying
space of a finite category with exactly two objects and aperiodic
endomorphism monoids. Its fundamental group is $\mathbb Z$, and
\[
 H_2(\widetilde X;\mathbb Z)\cong\mathbb Z[1/2].
\]
In particular, finite two-object categories are not restricted to wedges
of circles and simply connected summands, nor to finite-order monodromy.
\end{proposition}
\begin{proof}
Let $P$ be the five-element rectangular-band monoid realizing $S^2$.
Let $L$ be the monoid obtained by adjoining an identity to the
$3\times2$ rectangular band. Then
$BL\simeq S^2\vee S^2$ by Proposition~\ref{prop:B-band-wedge}.
The two unital embeddings $P\to L$ selecting row pairs $\{0,1\}$ and
$\{0,2\}$ induce the two integral Hurewicz generators, denoted $a,b$.
Apply the checked finite mapping-cone construction to a sphere map
representing $2a-b$. It gives a finite aperiodic monoid $M$ and a full
corner $eMe\cong L$, with
\[
 BM\simeq (S^2\vee S^2)\cup_{2a-b}D^3\simeq S^2,
\]
compatibly with the map from $BL$.
The last equivalence follows by a change of integral Hurewicz basis on
the wedge: the primitive vector $2a-b$ is one basis vector and its
attaching cell cancels that sphere. The images of $a,b$ in the remaining
$H_2\cong\mathbb Z$ are respectively one and two times a generator.

Compose the two embeddings with $L\cong eMe$. Lemma~\ref{D:corner-biset}
applied to $Me\amalg Me$ gives a finite two-object category whose
classifying space is the homotopy coequalizer of maps
$S^2\rightrightarrows S^2$ of degrees one and two. Replacing the
degree-one map by a homotopy equivalence identifies this with the
degree-two mapping torus. Its universal cover is the bi-infinite mapping
telescope of degree-two maps, so its second homology is the direct limit
of multiplication by two on $\mathbb Z$, namely $\mathbb Z[1/2]$.

Already the product of the two-parallel-arrow category with $P$ gives
a simpler two-object model of $S^1\times S^2$. The degree-two example
shows why eventual periodicity of an endomorphism of a fixed finite
monoid does not restrict all two-object realizations: the two maps here
have different source and target monoids.
\end{proof}

\subsection{Simultaneous finite realization of maps by full corners}

\begin{lemma}[A wedge construction preserving full corners]
\label{D:corner-wedge}
Let $P,Q$ be finite aperiodic monoids whose classifying spaces have finite
homotopy type. There is a finite aperiodic monoid $R$ and full-corner
embeddings of $P,Q$ into $R$ such that
\[
 BR\simeq BP\vee BQ
\]
under the two corner maps. Each left and right ideal joining a corner to
$1_R$ is a finite projective act over the corresponding corner monoid.
Here a finite projective act means a retract of a finite free act, or
equivalently a finite disjoint union of idempotent-generated principal
ideals.
\end{lemma}
\begin{proof}
The trivial factor can be omitted, retaining a minimal-ideal idempotent
for its one-element corner. Otherwise let
$K_P=I_P\times J_P$, $K_Q=I_Q\times J_Q$ be the minimal rectangular
bands, and choose $e_P\in K_P$, $e_Q\in K_Q$. Inside $P\times Q$ put
\[
 T=\{(1,1)\}\cup(P\times\{e_Q\})\cup(\{e_P\}\times Q)
       \cup(K_P\times K_Q).
\]
This is a submonoid: a product crossing the middle two pieces lies in
$K_P\times K_Q$. Its minimum is that rectangular band. The idempotents
$f=(1,e_Q)$ and $g=(e_P,1)$ satisfy
$fTf\cong P$, $gTg\cong Q$, because
$e_PK_Pe_P=\{e_P\}$ and $e_QK_Qe_Q=\{e_Q\}$.

Write $A_P=\mathbb ZP/\mathbb ZK_P$ and define $A_Q$ similarly. In
$A_T=\mathbb ZT/\mathbb Z(K_P\times K_Q)$ the idempotents
$f,g,h=1-f-g$ are central and orthogonal, giving
\[
 A_T\cong A_P\times A_Q\times\mathbb Z.
\]
For the augmentation modules of the row sets, the three summands are
\[
 fV_T\cong V_P,\qquad gV_T\cong V_Q,\qquad
 hV_T\cong V_P\otimes_{\mathbb Z}V_Q;
\]
the column augmentation modules satisfy the analogous identities.
These statements follow by writing the free row module as
$(\mathbb Z\oplus V_P)\otimes(\mathbb Z\oplus V_Q)$, using the chosen
rows as splittings of augmentation. The actions of $f,g$ project to the
first and second augmentation summands, while $h$ projects to their tensor
product.

Lemma~\ref{lem:B-band-Tor} and the product decomposition of $A_T$ now give
\[
 H_n(BT;\mathbb Z)\cong H_n(BP;\mathbb Z)\oplus H_n(BQ;\mathbb Z)
       \quad(n\ge3),
\]
and the same formula in degree two with an additional free summand of rank
\[
 d=(|I_P|-1)(|J_P|-1)(|I_Q|-1)(|J_Q|-1).
\]
The corner inclusions are split by the two coordinate projections
$T\to P,Q$; each opposite composite is constant at a minimal idempotent.
Thus their images form a split direct sum in homology. Since all groups
in each degree are finitely generated, the displayed abstract computation
shows that its complement is zero outside degree two and is
$\mathbb Z^d$ in degree two. All three classifying spaces are simply
connected, since their minimum ideals are bands. Choose Hurewicz spheres
for a basis of this complement. The resulting map
\[
 BP\vee BQ\vee\bigvee^d S^2\longrightarrow BT
\]
is an integral homology isomorphism between simply connected CW complexes,
hence a homotopy equivalence.

If $d>0$, use the finite mapping-cone construction to kill precisely these
$d$ sphere summands. It yields the required $R$ under $T$; if $d=0$, take
$R=T$. The embeddings remain full corners under this construction.

For completeness, before the cone attachment the right $P$-act $Tf$ is
\[
 P\ \amalg\!\coprod_{q\in Qe_Q\setminus\{e_Q\}} K_P.
\]
Every right $\mathcal R$-class of $K_P$ is an idempotent-generated right
ideal of $P$, so this act is projective. The analogous description proves
projectivity of $fT$, and likewise for $g$. In the explicit monoid
$M(Y,u)$ of Part B, the new left and right ideals at the old corner are
free acts over the old monoid: they are respectively $L\times[n]$ and
$Y_n\times L$. Iterating cone attachments consequently preserves
projectivity of all the older side acts, by the identities
$R f\cong R e\otimes_{eRe}(eRe)f$ and its left-hand counterpart.
This proves all assertions.
\end{proof}

\begin{lemma}[Relative cell trading]
\label{D:relative-cell-trading}
For a map $f:T\to X$ between finite connected simply connected CW
complexes, there is a finite relative CW complex $Z$ built from $T$ using
only cells of dimension at least two, and a homotopy equivalence $Z\simeq X$
whose restriction to $T$ represents $f$.
\end{lemma}
\begin{proof}
Replace $f$ by a cellular map and start with its finite mapping cylinder.
A relative maximal tree removes the relative zero-cells and their
connecting one-cells. Each remaining relative one-cell, completed to a
loop using paths in $T$, is null-homotopic in the mapping cylinder because
the latter has the homotopy type of $X$.

For every such one-cell add a two-cell along this completed loop and a
three-cell along the sphere formed by the new two-cell and a chosen null
disk in the old complex. These finitely many two/three-cell pairs leave
the homotopy type unchanged under $T$: the new two-cell adds a sphere,
and the new three-cell kills that sphere with coefficient one. Cellular
approximation places the chosen null disks in the old two-skeleton.
The subcomplex consisting of $T$, the relative one-cells, and the new
two-cells deformation retracts to $T$, since each new two-cell runs once
over its designated one-cell and otherwise lies in $T$. Replace this
subcomplex by $T$, transporting the remaining attaching maps along the
retraction. Homotopy invariance of pushouts along CW cofibrations preserves
the homotopy type under $T$. Only the old cells of dimension at least two
and the new three-cells remain.
\end{proof}

\begin{proposition}[Simultaneous full-corner realization]
\label{D:simultaneous-corners}
Let $T_1,\ldots,T_k,X$ be finite connected simply connected CW complexes,
and let $f_i:T_i\to X$ be a finite family of maps. Fix finite aperiodic
monoids $P_i$ realizing the $T_i$. There is a finite aperiodic monoid $M$
with full corners $e_iMe_i\cong P_i$ such that $BM\simeq X$ and the corner
maps represent the $f_i$. The side acts $Me_i$ and $e_iM$ are finite
projective $P_i$-acts.
\end{proposition}
\begin{proof}
Iterate Lemma~\ref{D:corner-wedge} to obtain a finite aperiodic monoid $P$
realizing $\bigvee_i T_i$, under full corners $P_i$, with projective side
acts. The maps $f_i$ assemble to a map from this wedge to $X$; basepoints
can be chosen compatibly because $X$ is simply connected. Apply
Lemma~\ref{D:relative-cell-trading}. Its finite sequence of cells of
dimension at least two can be attached successively using the checked
construction in Proposition~\ref{prop:B-finite-realization}: each attaching
sphere is connected, and each current monoid is aperiodic. Compatibility
under the old corner preserves the specified maps. At each step the old
corner is full and the new side acts are free over the old monoid, as
checked in the proof of Lemma~\ref{D:corner-wedge}. Composition therefore
preserves fullness and projectivity of the original $P_i$ corners.
\end{proof}

\begin{lemma}[A finite biset for a correspondence]
\label{D:projective-correspondence}
Suppose $P$ embeds as full corners $e_A M_A e_A$ and $e_B M_B e_B$, with
finite projective side acts. The finite $(M_B,M_A)$-biset
\[
 S=M_Be_B\otimes_P e_A M_A
\]
realizes the correspondence $BM_A\leftarrow BP\to BM_B$: its two-object
collage has the homotopy type of this homotopy pushout, compatibly with
the two vertex maps. A disjoint union of such bisets realizes the
corresponding disjoint union of edge spaces.
\end{lemma}
\begin{proof}
In the nerve of a two-object collage, a simplex crossing from the first
object to the second has a string of $M_A$ arrows, one element of $S$,
and a string of $M_B$ arrows. Splitting its barycentric coordinates at
the cross arrow gives the double mapping cylinder whose edge space is
the realization of this two-sided double bar construction.

The augmentation
\[
 B(M_Be_B,P,e_A M_A)\longrightarrow
 M_Be_B\otimes_P e_A M_A
\]
is a homotopy equivalence. Indeed it is so when its first right $P$-act
is free, by the extra degeneracy of the usual bar construction, and
finite projective acts are retracts of finite free acts. This equivalence
respects the outer actions. Substitute this bar resolution for $S$ in
the edge space and reassociate the bar directions. The two outer factors
become the action categories of the representable acts $M_Be_B$ and
$e_A M_A$; their classifying spaces are contractible, by the retract
argument of Lemma~\ref{D:corner-biset} and its dual. Replacing these
contractible factors by points leaves $B(*,P,*)=BP$. The two projections
are the corner maps, as follows from the same natural transformations
in the Karoubi categories used in Lemma~\ref{D:corner-biset}.
Realization and the bar reassociation preserve these comparisons because
all the simplicial spaces involved are CW spaces with cofibrant
degeneracy inclusions. Disjoint unions are preserved throughout.
\end{proof}

\begin{theorem}[Two objects realize finite graphs of simply connected spaces]
\label{D:graphs-two-objects}
Every finite connected graph of finite connected simply connected CW
spaces, with arbitrary continuous edge attaching maps, has the homotopy
type of the classifying space of a finite category with exactly two
objects and aperiodic endomorphism monoids.
\end{theorem}
\begin{proof}
Take the homotopy colimit over a maximal tree of the underlying graph.
It is a finite simply connected CW complex $Y$. The remaining edges
give finitely many pairs of maps $T_i\rightrightarrows Y$, with each
$T_i$ finite and simply connected. The original graph of spaces is the
homotopy pushout of two copies of $Y$ along the disjoint union
\[
 Y\amalg T_1\amalg\cdots\amalg T_r,
\]
where the first edge maps identically to both copies of $Y$, and the
other edge maps are the specified pairs. Collapsing the identity-edge
mapping cylinder recovers the original model.

Choose a finite aperiodic monoid $P_0$ for $Y$ and monoids $P_i$ for the
$T_i$. Apply Proposition~\ref{D:simultaneous-corners} separately to the
two families of maps into $Y$, including the identity map from $Y$.
This gives $M_A,M_B$ with the required common source corners and
projective side acts. Take the disjoint union of the finite bisets in
Lemma~\ref{D:projective-correspondence}, one for each of these $r+1$
edges. Their two-object collage realizes the required homotopy pushout.
\end{proof}

\begin{lemma}[A graph as the one-skeleton]
\label{D:rose-skeleton}
If $X$ is a finite connected CW complex with fundamental group $F_r$,
then $X$ has a finite CW model whose two-skeleton is
$R_r\vee\bigvee_{j=1}^a S^2$, where $R_r$ is a rose with $r$ petals.
\end{lemma}
\begin{proof}
Choose a map $R_r\to X$ inducing the given isomorphism on fundamental
groups, and take its finite mapping cylinder. The proof of
Lemma~\ref{D:relative-cell-trading} applies whenever the map on fundamental
groups is surjective: complete a relative one-cell to a loop, choose a
loop in the source representing the same element of the target fundamental
group, and use their difference as the attaching loop of the new two-cell.
This difference is null-homotopic in the target, so the compensating
three-cell exists. The new two-cell still traverses its designated
relative one-cell exactly once. Thus the same finite cell trading produces
a finite relative CW model over $R_r$ with no relative zero- or one-cells.

Each remaining two-cell attaches along a loop in $R_r$ that is null in
the total complex. Injectivity of $\pi_1(R_r)\to\pi_1(X)$ makes that loop
null-homotopic already in $R_r$. Replacing these finitely many attaching
maps by constant maps, and transporting the higher attaching maps through
the resulting equivalence of two-skeleta, gives the asserted model.
\end{proof}

\begin{lemma}[Enlarging pieces over a tree]
\label{D:tree-enlargement}
Let $Z$ be a finite graph of finite connected simply connected CW spaces
over the rose $R_r$, and let $n\geq3$. For any map
$\alpha:S^{n-1}\to Z$, the space $Z\cup_\alpha D^n$ again has a model of
this form over the same rose.
\end{lemma}
\begin{proof}
We may assume $r\geq1$ and that all attaching maps are cellular. Write
$G=F_r$ and let $T$ be the Cayley tree for its displayed free basis.
The universal cover $Y$ of $Z$ is the tree of the prescribed simply
connected vertex and edge spaces. It has its usual projection $q:Y\to T$:
vertex spaces map to vertices, and edge mapping cylinders map to edges.
For every nonempty finite subtree $K\subset T$, the inverse image
$q^{-1}(K)$ is a finite simply connected CW complex. This follows by
successively adjoining leaf vertex spaces along connected simply
connected edge spaces, using van Kampen; it also follows directly from
the graph-of-spaces homotopy colimit.

Lift $\alpha$ to $\widetilde\alpha:S^{n-1}\to Y$, and choose a finite
subtree $S\subset T$ containing $q\widetilde\alpha(S^{n-1})$. Compactness
and local finiteness of $T$ permit this choice. For an integer $R\geq1$
large enough that some combinatorial radius-$R$ ball contains $S$, put
\[
 D_g=q^{-1}(B_R(g)) \qquad(g\in G).
\]
Here $B_R(g)$ is the full subtree on vertices at distance at most $R$
from $g$. For every edge $[g,h]$ of $T$, put
$E_{gh}=D_g\cap D_h=q^{-1}(B_R(g)\cap B_R(h))$.
Both types of pieces are finite and simply connected, because the
intersection of these two balls is a nonempty finite subtree.

The tree homotopy colimit with vertex pieces $D_g$ and edge pieces
$E_{gh}$ maps naturally to $Y$, and this map is a homotopy equivalence.
Here is a cellular justification that also keeps track of the finite
quotient. For each open cell $\sigma$ of the original tree of spaces,
its set of membership vertices
\[
 A_\sigma=\{g:\sigma\subset D_g\}
\]
is a nonempty finite subtree of the indexing tree $T$. Indeed $q(\sigma)$
is either a vertex or an open edge, and the membership set is respectively
one ball of centers or the intersection of the two balls of centers
determined by its endpoints. At the cellular stage for $\sigma$, the
homotopy colimit therefore adjoins
$D^{\dim\sigma}\times |A_\sigma|$ along
$S^{\dim\sigma-1}\times |A_\sigma|$; the projection to the corresponding
cell of $Y$ is a homotopy equivalence of these attaching pairs because
$|A_\sigma|$ is contractible. Induction over the skeleta and the CW gluing
lemma prove the assertion. The construction is $G$-equivariant, with no
cell stabilizers. The same cellular argument on cell orbits proves that
the quotient graph homotopy colimit is homotopy equivalent to $Z$.

For each translated attaching map $\gamma\widetilde\alpha$, define
\[
 A_\gamma=\{g:\gamma S\subset B_R(g)\}
          =\bigcap_{v\in\operatorname{Vert}(\gamma S)}B_R(v).
\]
This is a nonempty finite subtree: it is nonempty by the choice of $R$
and translation, and intersections of convex subtrees of a tree are
convex. Attach a copy of its $n$-cell to every $D_g$ for $g\in A_\gamma$,
and attach the same cell to $E_{gh}$ whenever $[g,h]\subset A_\gamma$.
The boundary map lies in all the required pieces by the definition of
$A_\gamma$. Call the enlarged vertex and edge pieces $D'_g,E'_{gh}$.
They are simply connected since only cells of dimension at least three
have been added. They are finite: for a fixed $g$, only finitely many
translates $\gamma S$ lie in $B_R(g)$, since the action on vertices is
free and the ball is finite.

The new tree homotopy colimit is obtained from the old one by adjoining,
for each $\gamma$, the pair
\[
 (D^n\times |A_\gamma|,
  S^{n-1}\times |A_\gamma|).
\]
Its natural map to $Y$ with all translates of the new cell attached
collapses the contractible tree $|A_\gamma|$. Homotopy invariance of the
attaching pushouts, or the same cellular argument, makes this map a
homotopy equivalence. Everything is $G$-equivariant. Taking the quotient
gives a finite graph of the simply connected spaces $D'_1$ and
$E'_{1,s}$, one edge for each free generator $s$, and a homotopy
equivalence from its total space to $Z\cup_\alpha D^n$.
\end{proof}

\begin{theorem}[Finite complexes with free fundamental group]
\label{D:free-pi-two-objects}
A finite connected CW complex $X$ is homotopy equivalent to the
classifying space of a finite category with exactly two objects if and
only if its fundamental group is free. In the sufficient direction the
endomorphism monoids can both be chosen aperiodic.
\end{theorem}
\begin{proof}
For sufficiency, Lemma~\ref{D:rose-skeleton} starts a finite CW model of
$X$ with $R_r\vee\bigvee_a S^2$. This is a graph of simply connected
finite spaces over $R_r$: use $\bigvee_a S^2$ as the vertex space and a
point for each edge space. Attach all the finitely many higher cells,
one at a time, using Lemma~\ref{D:tree-enlargement}. The resulting finite
graph of finite simply connected spaces is realized by the two-object
category of Theorem~\ref{D:graphs-two-objects}.

For necessity, Theorem~\ref{A:bounded-categories} kills every endomorphism
loop in the fundamental groupoid of any finite category with finite
CW type. If both cross-hom sets of a two-object category are nonempty,
choose arrows $a:0\to1$ and $b:1\to0$. Their composites are endomorphisms,
so $b=a^{-1}$ in the fundamental groupoid; any other cross arrow equals
$a$ or $b$ by the same argument. The fundamental group is then trivial.
Otherwise the connected category is triangular. After its endomorphisms
have been killed, its fundamental-group presentation has one generator
for each equivalence class of cross arrows under left and right
endomorphism actions and no other relations. Choosing one such class
as a spanning-tree edge leaves a free group on the remaining classes.
\end{proof}

\begin{corollary}[The full two-object classification]
\label{D:two-object-remaining}
The connected finite-dimensional CW homotopy types realized by finite
categories with exactly two objects are precisely the finite connected
CW homotopy types with free fundamental group.
\end{corollary}
\begin{proof}
Theorem~\ref{D:two-objects} proves that a finite two-object category with
finite-dimensional classifying-space type has finite CW type and free
fundamental group: finite domination, vanishing of
$\widetilde K_0(\mathbb ZF_r)$, and Wall's theorem supply finiteness.
The converse is Theorem~\ref{D:free-pi-two-objects}.
\end{proof}

\begin{remark}\label{D:tree-repair-scope}
The tree enlargement proof uses no homology-preserving plus construction
on the individual vertex or edge pieces. Their homology is allowed to
change; the product with a contractible membership subtree controls the
homotopy type of the total space.
\end{remark}

\begin{corollary}\label{D:two-object-disconnected}
The disconnected finite-dimensional types realized by categories with
exactly two objects are precisely disjoint unions of two nonempty
finite connected simply connected CW homotopy types.
\end{corollary}
\begin{proof}
A disconnected two-object category has no cross arrows and is the
disjoint union of its two endomorphism monoids. Section~\ref{sec:A}
forces each finite-dimensional component to be simply connected and
of finite CW type. Conversely Theorem C, checked in Section~\ref{sec:B},
realizes the two given simply connected finite types separately by
finite monoids; their disjoint union is the required category.
\end{proof}
